% 本文件由 LaTeX 排版 Agent 根据有效 translation.json 自动生成。
% author=Tertullian; work_id=0304; title=论冠冕

\chapter{\WorkTitle{论冠冕}}

% structure: p0001 source=h1 target=omitted reason=page-title-already-rendered-by-parent

\Emph{论冠冕，或称《论冠冕》。}\par

% structure: p0003 source=h2 target=section reason=relative-heading-rank; base=h2

\section{\Emph{第一章}}

最近发生了一件事：正当我们最优秀的皇帝们的恩赐在军营中分发时，士兵们头戴桂冠走上前来。其中一人，更是天主的士兵，比其余那些自以为能事奉两个主人的弟兄们更加坚定，他独自光着头，手中拿着无用的冠冕——甚至凭此特点已被众人认出是基督徒——高贵的引人注目。于是，所有人都开始指认他，远处嘲笑，近处咬牙切齿。那人在刚离开队列时，低语声传到了指挥官那里。指挥官立刻问他：\Quote{你的装束为何如此与众不同？}他宣称自己无权与其他人一样戴冠冕。被催问原因时，他回答：\Quote{我是基督徒。}啊，士兵！你在天主内夸耀。于是案件被审议并表决；事情被提交到更高的法庭；犯事者被带到长官面前。他立即脱下厚重的斗篷，开始卸下重负；他从脚上松解军鞋，开始站在圣地上；他放下宝剑，这剑对于保护我们的主并非必要；从手中也掉落了桂冠；如今，身披自己鲜血希望的紫色，脚穿福音预备的鞋子，腰束天主更锐利的话语，全副武装着宗徒的甲胄，头戴更为相称的殉道白冠，他在监狱中等待着基督的赏赐。此后，对他的行为开始有不利的评判——我不知道是否来自基督徒，因为外教人的评判并无不同——仿佛他固执鲁莽，过于急于赴死，因为仅仅因着装一事被责问，便给那些尊名者带来麻烦——他，确实，在如此众多的士兵弟兄中唯独勇敢，唯独是基督徒。显然，正如他们拒绝了圣神的预言，他们也意图拒绝殉道。于是他们抱怨说如此美好长久的和平因他们而处于危险。我也不怀疑有些人已经在背弃圣经，收拾行囊，准备好从一城逃往另一城；因为他们只记得福音中这些内容。我也知道，他们的牧者在和平时是狮子，争战时却是鹿。至于那些为逼取我们认罪而提出的问题，我们将在别处教导。现在，既然他们也提出异议——\Quote{但我们何处被禁止戴冠冕？}——我将就此点加以论述，这更适合在此处理，事实上它是当前争论的核心。这样，一方面，那些无知但渴望求知的研究者可以受到指导；另一方面，那些试图为罪辩护的人，尤其是戴桂冠的基督徒自己，他们把这仅仅当作辩论的问题，好像它要么根本不是过犯，要么至少是可疑的，因为可以加以调查。然而，我现在将证明它既不是无罪的，也不是可疑的。\par

% structure: p0005 source=h2 target=section reason=relative-heading-rank; base=h2

\section{\Emph{第二章}}

我断言，没有一个信友头上戴过冠冕，除非在受试探的时候。这是所有人的情况，从慕道者到宣信者、殉道者，或（视情况而定）背教者。那么，请思考一下我们现在主要探讨的这个习惯从何处获得权威。但当有人问及为何遵守它时，同时清楚的是它被遵守。因此，违反一个可以凭其名声辩护、且得到普遍接受支持的习惯，既不能被视为无罪，也不能被视为不确定。它是毋庸置疑的，所以我们应当探究此事的原因；但不应损害习惯，目的不是推翻它，而是建立它，以便你在对原因满意后，能更加谨慎地遵守它。但是，一个人违反习惯后，又将习惯拿出来辩论，并且因为停止遵守而寻求曾经拥有它的解释，这是何等做法？因为，尽管他可能因此希望显得渴望探究，以表明他放弃它没有做错，但显然他之前已经因为擅自遵守而犯了过犯。如果他今天接受冠冕没有做错，那么他先前拒绝它就已经犯了过犯。因此，这篇论文不是为了那些不适宜探究的人，而是为了那些怀着真正受教愿望的人，他们提出的不是辩论的问题，而是请求建议。因为真正的探究总是出于这种愿望；我赞扬那信德，它相信在了解原因之前遵守规则是当尽的本分。立即质问何处写着我们不应戴冠冕是容易的。但哪里写着我们应戴冠冕呢？事实上，当人们急迫地要求圣经支持他们不同的一方时，他们预先判断圣经的支持同样应该出现在他们那一边。因为如果有人说戴冠冕是合法的，理由是圣经没有禁止它，那么同样有效地可以反驳说，正是因为这个理由戴冠冕是非法的，因为圣经没有命令它。纪律该怎么办？是否接受两者，好像两者都没有被禁止？或者拒绝两者，好像两者都没有被命令？但是\Quote{没有被禁止的事是自由允许的。}我宁愿说，没有被自由允许的事是被禁止的。\par

% structure: p0007 source=h2 target=section reason=relative-heading-rank; base=h2

\section{\Emph{第三章}}

然而，当我们拥有一项古老惯例，它已预先为我们确立了问题所处的状态，即问题的本质时，我们还要在这条线上来回拉锯多久呢？如果圣经中没有章节规定它，那么无疑，源自传统的习惯已确认了它。因为若未经传授，任何事怎能成为惯例呢？即使在引证传统时，你说也必须要求书面权威。因此，让我们探究：传统若非成文，是否不应被接纳？当然，我们会说，若没有其他惯例（我们仅凭传统及其后习惯的支持来维持，无需任何书面文件）的先例，它就不应被接纳。为简明起见，我从圣洗圣事开始。当我们即将进入水中时，稍早之前，在会众面前，在主持者手下，我们庄重地宣誓弃绝魔鬼及其虚荣及其天使。随后我们三次浸入水中，所作的承诺比主在福音中所指定的更为充分。然后当我们被扶起（如同新生婴儿）时，我们首先品尝牛奶和蜂蜜的混合，并从那天起一周内不进行日常沐浴。我们也在天亮前、且仅从主持者手中，在聚会中领受圣体圣事——主曾命令在用餐时间领受，并吩咐众人都领受。每当周年纪念来临，我们为亡者献祭，作为生日荣誉。我们视在主日禁食或跪拜敬拜为不合法。从复活节到圣神降临节，我们也享有同样特权。若有任何酒或面包（即便是我们自己的）掉落地上，我们感到痛心。在每一步前进和移动，每一次进出，穿衣穿鞋时，沐浴时，就座进餐时，点灯时，卧床、坐席，以及日常生活的所有普通行动中，我们在额上画十字圣号。\par

% structure: p0009 source=h2 target=section reason=relative-heading-rank; base=h2

\section{第四章}

若你坚持要为这些及其他类似规则找到明确的圣经命令，你将一无所获。传统会被展示为它们的创始者，习惯为它们的强化者，信德为它们的遵守者。理性会支持传统、习惯和信德，这一点你或自己领悟，或从他人了解。同时，你应相信存在某种应当服从的理由。我再加一个实例，因为适当展示古人的做法也有益。在犹太人中间，妇女蒙头如此普遍，以致可借此认出她们。我在此要求律法。我将宗徒暂放一旁。倘若\Person{黎贝加}远远看见她的未婚夫时就拉下帕子，这区区个人的谦逊不可能成为律法，或只对具有她那样理由的人成立。让贞女们独自蒙头吧，且仅在即将结婚、认出未来丈夫之时。若还有\Person{苏撒纳}，她在受审时被掀开面纱，为妇女蒙头提供了论据，我在此也可以说，面纱是自愿的。她是被控告而来，为自己带来的耻辱而羞愧，恰当地隐藏美貌，甚至因现在害怕取悦他人。但我不认为，当她的目的是取悦时，她会在丈夫的林荫道上蒙头散步。姑且假设她总是蒙头。在此特定情况，或事实上在任何其他情况，我要求服饰律法。若我到处找不到律法，那么传统已将这种时尚赋予了习惯，并后来在宗徒的认可中（从理性的真正解释）找到权威。因此，这些实例足以说明：你可以辩护甚至维持由习惯确立的未成文传统；传统通过长期遵守而得到适当见证。但在民事事务中，当缺乏成文法律时，习惯也被接受为法律；无论它基于文字还是理性，结果相同，因为理性实为法律的基础。但你说，若理性是法律的基础，那么从现在起，任何人提出以理性为基础的东西，都将被视为法律。或者，你认为每个信徒都有权创立和建立法律，只要它合乎天主、有助于纪律、促进救恩，而主说：\Quote{但你们为什么连自己判断什么是正义的事也不呢？}\BibleRef{路 12:27}？不仅就司法判决，而且就我们被召思考的每一决定，宗徒也说：\Quote{若你们对什么事无知，天主会启示给你们；}\BibleRef{斐 3:15}他本人也惯于提供建议，尽管没有主的命令，并凭自己发言，因为他拥有引导进入全部真理的天主之神。因此，他的建议凭借神圣理性的确证，已等同于神圣命令。现在，请认真向这位师傅求教，同时保持你对传统的敬意，无论它最初源于何人；不要顾及作者，而要看重权威，尤其是习惯本身的权威——正因如此，我们应当尊敬它，以免缺少解释者；这样，若理性也是天主的恩赐，你便可学会：并非习惯是否应由你遵循，而是为什么。\par

% structure: p0011 source=h2 target=section reason=relative-heading-rank; base=h2

\section{第五章}

基督徒实践的理由更加有力，因为自然（它是万物的第一个法则）也支持它们。自然是最先断定花冠不配于头的。不过，我认为我们的天主是自然的天主，祂创造了人；并且，为了使人渴望（欣赏、分享）祂所造之物带来的欢愉，祂赋予人某些感官，通过这些感官的（运作的）肢体，可以说，这些肢体是它们的专属工具。听觉安置在耳朵里；视觉点亮在眼睛中；味觉封闭在口内；嗅觉飘入鼻子；触觉固定在指尖。通过这些外在之人的器官为内在之人服务，神圣恩赐的享受经由感官传递给灵魂。那么，花朵中有什么给你带来享受？因为田间的花朵正是花冠的特殊（至少是主要）材料。你说，要么是香气，要么是颜色，或者两者兼有。颜色和香气的感官是什么？我想是视觉和嗅觉。这些感官被分配给了哪些肢体？如果我没有弄错的话，是眼睛和鼻子。那么，就用视觉和嗅觉来享用花朵吧，因为正是这些感官被指定来享受它们；通过属于这些感官的肢体，即眼睛和鼻子来使用它们。你从天主得到了事物本身，从世界得到了方式；但非凡的方式并不妨碍以通常的方式使用事物。那么，让花朵——无论是彼此系在一起，还是用线和灯心草绑在一起——保持它们自由、松散时的样子，成为被观看和嗅闻的东西。比如，当你把一束花串成一串时，你称之为花冠，这样你可以一次携带许多，同时享受它们。那么，如果它们极其纯净，就放在你的怀里；如果它们极其柔软，就撒在你的床上；如果它们完全无害，就放进你的杯中。以尽可能多的、迎合你感官的方式享受它们。但是，头——它既不能分辨颜色，也不能吸入芬芳，也不能感受柔软——对花朵有什么品味？对花冠的任何部分（除了其带子）有什么感官呢？用头去渴求花朵，与用耳朵去渴求食物，或用鼻孔去渴求声音，同样是违反自然的。然而，凡是违反自然的事，都理应被所有人视为怪异；但对我们来说，它还要被谴责为对天主——自然的主和创造者——的亵渎。\par

% structure: p0013 source=h2 target=section reason=relative-heading-rank; base=h2

\section{第六章}

那么，既然要求天主的法律，你就有一条普世通行的法律，刻在自然的心版上，宗徒也惯常引用它，就如论到妇女的帕头时他说：\Quote{连自然本身没有教导你们吗？}\BibleRef{格前 11:14}——就如写信给罗马人时，肯定外邦人按本性遵行法律所要求的事，\BibleRef{罗 2:14}，他暗示了自然律和彰显律法的自然。是的，在同一书信的第一章，他也确证了自然，当他断言男女之间彼此改变了受造物自然的用途，成为不自然的，\BibleRef{罗 1:26}，作为他们错误的刑罚报应。我们最初确实藉着自然的教导认识天主本身，称祂为万神之神，假定祂是善的，并呼求祂为审判者。难道你还有疑问：为了享受祂的受造物，自然是否应是我们的向导，以免我们被天主的对手引向歧途？那对手已腐化了人类本身，也腐化了为我们的种族指定用途的整个受造界，因此宗徒说受造物也不甘情愿地屈服于虚空，完全丧失了原有的特质，先是被虚浮，后是被卑鄙、不义、不敬的使用所败坏。同样，在娱乐的享乐中，受造物被那些凭本性确实感知到所有组成娱乐的材料都属于天主，却缺乏认知去察觉它们已被魔鬼改变的人所羞辱。但是，为了我们自己的戏剧爱好者，我们已经充分讨论了这个主题，而且是用希腊文写成的著作。\par

% structure: p0015 source=h2 target=section reason=relative-heading-rank; base=h2

\section{第七章}

让这些冠冕商人在此期间承认自然的权威——基于作为人类共有的常识——以及他们特殊宗教的证明，即按上一章所说，他们是自然的天主的崇拜者；并且，好似超出所需，让他们也考虑那些禁止我们戴冠冕的其他理由，尤其是戴在头上的冠冕，实际上是一切种类的冠冕。因为我们不得不偏离我们与人类共享的自然法则，为保持我们基督徒规范的全部特性，这也关系到其他种类的冠冕，它们似乎由不同物质组成，为不同用途而设；否则，因为它们不是由花制成——而花的用途是自然所指示的（正如在这个军事月桂冠的情况中所做的那样），它们可能被认为不受我们教派的禁令约束，因为它们避开了自然的任何反对。那么，我看我们必须更深入、更充分地探讨这个问题，从其起源开始，经过其连续的发展阶段，直到其更不规则的演变。为此，我们需要转向异教文献，因为属于异教的事物必须从他们自己的文献中证明。我所得的这一点点，我相信就足够了。如果真有一个\Person{潘多拉}，\Person{赫西俄德}提到她是第一个女人，那么她的头是第一个被美惠女神加冕的，因为她从所有\Emph{神}那里得到礼物，因而得了\Emph{她的名字}潘多拉。但是\Person{梅瑟}，一位先知，而非牧羊诗人，给我们显示第一个女人\Person{厄娃}，她的腰间更自然地围上叶子，而非她的太阳穴戴上花朵。那么，潘多拉是一个神话。因此，我们不得不为冠冕的起源而羞愧，即使是基于与之相关的谬误；并且，正如很快将要显现的，同样也基于其现实。因为无疑，某些人要么开创了此物，要么为其增添了光彩。\Person{费雷基德斯}叙述说，\Person{萨图尔努斯}是第一个戴冠冕的人；\Person{狄奥多罗斯}说，\Person{朱庇特}在征服\Person{泰坦}后，受到其他神祇的这份礼物致敬。同一位作者也将头带分配给\Person{普里阿普斯}；给\Person{阿里阿德涅}一个金和印度宝石的花环，是\Person{伏尔甘}的礼物，后来归\Person{巴克斯}，随后变成了一个星座。\Person{卡利马科斯}给\Person{朱诺}戴上了葡萄藤冠冕。同样在\Place{阿尔戈斯}，她的雕像戴着葡萄藤花环，脚下放置一张狮皮，展示这位继母因从两个继子那里获得的战利品而欢欣。\Person{赫丘利}头上时而戴白杨，时而戴野橄榄，时而戴欧芹。你有关于\Person{刻耳柏洛斯}的悲剧；你有\Person{品达}；还有\Person{卡利马科斯}，他提到\Person{阿波罗}在杀死德尔斐之蛇后，作为一个恳求者，戴上了月桂花环；因为在古代，恳求者通常被加冕。\Person{哈波克拉提翁}论证说，\Person{巴克斯}（与埃及人中的\Person{奥西里斯}相同）故意被戴上常春藤冠冕，因为常春藤的性质是保护大脑免受困倦。但另一方面，\Person{巴克斯}也是月桂花冠（冠冕）的创始者，他戴着它庆祝他对印度人的胜利，就连普通民众也承认，他们称奉献给他的日子为\Quote{大冠冕}。如果你再翻开埃及人\Person{利奥}的著作，你会学到\Person{伊西斯}是第一个发现并将麦穗戴在头上的人——这是更适合肚子的东西。那些想要更多信息的人，将会在\Person{克劳狄乌斯·萨图尔尼努斯}（一位才华出众的作家，也处理这个问题）的著作中找到关于此主题的充分阐述，因为他有一本关于冠冕的书，如此解释了它们的起源、原因、种类和仪式，以至于你会发现在花中所有迷人的东西、在叶枝中所有美丽的东西，以及每一块草皮或葡萄藤嫩枝都已奉献给某个人头；这充分说明我们应当判断戴冠的习俗对我们有多么陌生，这个习俗是由那些世人相信为神祇的人引入，并此后不断为他们荣耀而管理的。如果魔鬼，从起初就是说谎者，甚至在这件事上为他虚假的神性体系（偶像崇拜）工作，那么他自己也无疑为他的神性谎言被实现提供了条件。那么，被万民引入以尊崇魔鬼候选人的东西，从一开始就专为这些而设，且甚至在那时通过偶像和在仍然活着的偶像中接受了对偶像崇拜的祝圣，这东西在真天主子民中当算作什么？并非仿佛偶像是什么，而是因为别人献给偶像的东西属于魔鬼。但如果\Emph{别人}献给偶像的东西属于魔鬼，那么当偶像活着时它们自己献给自己的东西，岂不更是如此！魔鬼自己，无疑，在欲望和渴求的状态中，通过他们所附身的人，在实际供应之前已为自己做了准备。\par

% structure: p0017 source=h2 target=section reason=relative-heading-rank; base=h2

\section{第八章}

暂且持守这种信念，当我审视一个挡在我们面前的问题时。因为我已听见有人说，许多其他事物，如同冠冕一样，皆由世人信为神明的那些人所发明，然而，它们仍见于我们现今的习惯、早期圣人的习惯、对天主的敬礼，以及\Person{基督}自身身上——祂以人的身份行事时，所使用的无非就是这些普通的人生工具。好吧，就如此吧；我也不再进一步追究这些事物的起源。假设\Person{墨丘利}是第一个教授文字知识的人；我承认，文字对于生活的经营和交易，以及履行我们对天主的虔诚，都是必需的。不仅如此，即便祂是第一个弹拨琴弦发出旋律的人，当我聆听\Person{达味}时，我不会否认这一发明已被圣人们使用，并服务于天主。假设\Person{阿斯克勒庇俄斯}是第一个寻求并发现疗法的人：\Person{依撒意亚}\BibleRef{依 38:21}提到，当他生病时，他命令希则克雅用药。\Person{保禄}也知道少许酒对胃有益。\BibleRef{弟前 5:23}假设\Person{密涅瓦}是第一个建造船只的人：我将看见\Person{约纳}和宗徒们航行。不止这些：因为甚至\Person{基督}，我们将会发现，也穿着普通的衣服；\Person{保禄}也有他的外衣。如果你立刻将每一件家具和每一个家用器皿，都归功于某个世间神明的发明，那么，我必须承认\Person{基督}，无论是祂斜靠在躺椅上、为门徒洗脚端来盆、从水罐中倒水，还是用一条细亚麻毛巾束腰——这是一件特别供奉给\Person{奥西里斯}的服饰\BibleRef{若 13:1}。大体上我如此回应这一点，确实承认我们与其他人一起使用这些物品，但要求根据合理与不合理的事物之间的区别来判断这一点，因为无区分地使用它们具有欺骗性，隐藏了受造物的败坏，借此它已屈服于虚荣。因为我们主张，只有那些满足人类生活必需品、提供单纯有用之物、给予真正帮助和体面舒适的东西，才适合我们和前人使用，才唯独适合对天主和\Person{基督}本身的服侍。因此，它们可以合理地被认为是来自天主自己的灵感，天主无疑首先为祂自己的人（\OriginalTerm{人类}{man}）的享受提供了、教导了并服侍了。至于不属于此类的事物，它们不适合在我们中使用，尤其是那些因此既不在世界中、也不在基督的道路上所发现的东西。\par

% structure: p0019 source=h2 target=section reason=relative-heading-rank; base=h2

\section{第九章}

简而言之，你曾发现哪一位圣祖、先知、肋未人、司铎、统治者，或后来的哪一位宗徒、福音宣讲者、主教，是冠冕的佩戴者？我想连天主的圣殿本身也没有戴冠；同样，约柜、会幕、祭台、灯台也没有戴冠，尽管确实，在第一次奉献的隆重典礼和第二次重建的欢庆中，戴冠本应是最合适的，如果这对于天主是相称的话。但如果这些事物是我们（因为我们是天主的殿、祭台、光明和圣器）的预像，那么它们也以预像表明：天主的子民不应戴冠。现实必须始终与形象相符。如果你可能反驳说\Person{基督}本人曾被戴冠，那么你将得到简短的答复：你也像祂那样戴冠吧；你完全有自由。然而，就连那狂妄不敬的冠冕，也并非来自犹太人的任何法令。那是罗马士兵的发明，取自世俗的习俗——这种习俗天主的子民从未允许，无论是在公共欢庆的日子还是为了满足天生的奢侈：所以他们从巴比伦俘囚中归来时，带着手鼓、笛子和瑟，比带着冠冕更合适；他们吃喝之后，不戴冠，起来玩乐。关于这欢乐的记述和奢侈的揭露，也都不会对冠冕的荣耀或羞辱保持沉默。同样，\Person{依撒意亚}说：\Quote{用手鼓、瑟和笛饮酒}\BibleRef{依 5:12}，如果这种习俗曾在属于天主的事物中占有一席之地，他本会加上\Quote{戴着冠冕}。\par

% structure: p0021 source=h2 target=section reason=relative-heading-rank; base=h2

\section{第十章}

因此，当你声称异教神祇的装饰品也出现在天主那里，目的是要从中声称头冠可以普遍使用时，你早已为自己定下规则：凡是在事奉天主中找不到的东西，我们就不能把它当作可与他人共享使用之物而保留在我们中间。\newline{}那么，有什么比配得上偶像的东西更不配天主的呢？又有什么比配得上死人的东西更配得上偶像的呢？因为死人也享有戴上这样的冠冕的特权，因为他们立即借着他们的服饰和神化的敬礼变成了偶像，这种神化在我们看来是第二次的偶像崇拜。\newline{}既然缺乏感觉，他们就拥有了使用那缺乏感觉之物的权利，就好像他们完全拥有感觉却想要滥用它一样。当使用中不再有任何实在性时，使用和滥用就没有区别了。一个人想要实现其目的的感知本性并不为他所用，他怎能滥用某物呢？\newline{}此外，宗徒禁止我们滥用，而他本会更自然地教导我们不要使用，除非基于这样的理由：凡是对事物毫无感知的地方，也就没有对它们的错误使用。\newline{}但整个事情毫无意义，就偶像而言，事实上是一件死的工作；然而，毫无疑问，就属于宗教礼仪的魔鬼而言，它是一件活的工作。\newline{}\BibleQuote{外邦人的偶像，}\Person{达味}说，\BibleQuote{是金银。}\BibleQuote{他们有眼睛，却看不见；有鼻子，却闻不到；有手，却不能触摸。}诚然，我们本应通过这些器官来享受鲜花；但若他声明制造偶像的人将会像它们一样，那么凡是以偶像装饰的风格使用任何东西的人，已然如此了。\newline{}\BibleQuote{洁净的人，一切都洁净；同样，不洁净的人，一切都不洁净；}\BibleRef{铎 1:15}但没有什么比偶像更不洁净的了。\newline{}这些物质本身作为天主的受造物并无不洁，在其本然状态下，人人皆可使用；但它们在用途上所投入的礼仪造成了所有差异；因为我自己也杀一只鸡为自己，就像\Person{苏格拉底}为\OriginalTerm{阿斯克勒庇俄斯}{Æsculapius}所做的那样；如果某个地方的臭味使我厌恶，我也自己焚烧阿拉伯产品，但并非以与偶像崇拜同样的仪式、同样的服饰、同样的排场来行。\newline{}如果受造物仅仅因一句话就变为不洁，如宗徒所教导的：\BibleQuote{但若有人说：这是祭献给偶像的，你们就不要碰它，}\BibleRef{格前 10:28}那么当它被祭献给诸神的服饰、礼仪和排场污染时，更是如此。因此，冠冕也被证明是献给偶像的供物；因为借着这仪式、服饰和排场，它被祭献给偶像，而偶像是它的创始者，其使用特别归属它们；主要基于这个原因：凡是在天主的事物中没有地位的，就不可以被我们像他人一样接纳使用。\newline{}因此宗徒呼喊：\BibleQuote{逃避偶像崇拜！}\BibleRef{格前 10:14}他所指的无疑是全部的完整的偶像崇拜。\newline{}请反省那是什么样的丛林，其中隐藏着多少荆棘。不可给偶像任何东西，因此也不可从偶像那里取任何东西。如果在偶像庙宇中斜倚与信德不符，那么穿着偶像的服饰又算什么？基督与\OriginalTerm{彼列}{Belial}有什么相合？因此，要逃避它；因为他命令我们远离偶像崇拜——不与它有任何密切来往。甚至地上的蛇也能用它的气息将人吸入远处。更进一步，若望说：\BibleQuote{我的孩子们，你们要远避偶像，}\BibleRef{若一 5:21}——现在不是远离偶像崇拜，仿佛远离其敬礼，而是远离偶像本身——即远离任何与它们相像的事物：因为你们这些永生天主的肖像，竟成为偶像和死人的模样，这是不相称的事。\newline{}到目前为止，我们断言，这种服饰属于偶像，无论是从其起源历史来看，还是从假宗教的使用来看；此外，基于这个理由：虽然它没有被提及与敬拜天主相关，但它越来越多地被交给那些在其古迹、节庆和礼仪中发现它的人。简言之，连门、祭品和祭台、仆人和司祭，都戴着冠冕。你在\Person{克劳狄}那里可以看到各种司祭团体的冠冕。\newline{}我们还增加了完全不同的事物之间的区别——即合理的事物与违背理性的事物——以回应那些因为某些东西有共同使用而主张参与一切事物之权利的人。\newline{}因此，关于这部分主题，现在仍需检查戴冠冕的特殊理由，当我们证明这些理由是陌生的，甚至违背我们的基督徒规诫时，我们就能证明它们没有任何理性的辩护依据，借此这种服饰可以被辩护为我们可参与使用之物，正如甚至某些其他人所举的例子那样。\par

% structure: p0023 source=h2 target=section reason=relative-heading-rank; base=h2

\section{第十一章}

首先，关于军冠的真正根据，我认为我们应先探问，战争对基督徒而言是否全然合适。既然其基础已被定罪，讨论其偶然因素又有何意义？我们是否认为在神圣誓言之外再加上人的誓言是合法的？即在基督之后，人又向另一位主人承诺，并弃绝父母和所有至亲——甚至法律也命令我们仅次于天主而尊敬和爱他们，而福音也同样、只认为他们仅次于基督而给予尊敬？主宣告凡动刀的必死于刀下，又怎能认为以刀剑为业是合法的呢？和平之子岂能参与战斗，何况他连诉讼也不适宜？他连自己的冤屈也不报复，又怎能施加锁链、监禁、酷刑和惩罚呢？\Emph{他}难道要替别人守夜多于替基督？或要在主日守夜，而他甚至不为基督自己守夜？他岂能在已弃绝的庙宇前站岗？他岂能在宗徒禁止的地方就食？\BibleRef{格前 8:10}他岂能用刺穿基督肋旁的长矛靠着歇息，夜间殷勤保护那些他白天通过驱魔赶走的人？他岂能举着敌对基督的旗帜？\Emph{他}已从天主那里领受了口令，岂能再向皇帝索要口令？\Emph{他}期待被天使的号角唤醒，岂能因吹号者的号角在死亡中受惊扰？基督徒不被允许向偶像烧香，基督已免除他火的刑罚，岂能按照营规被焚烧？此外，营中职务所涉的众多过犯——我们必须视其为违犯天主的律法——只需略加审视便可看到。将这名号从光明之营带到黑暗之营本身就是违犯。当然，若信德后来才来到，发现某人已投身军旅，其情况则不同，如同若翰曾为之施洗的那些人，以及那些最忠信百夫长——我是指基督赞许的那位百夫长和伯多禄教诲的那位百夫长。但同时，当一个人成为信徒，信德已被封印，就必须要么立即放弃军职——这是许多人的做法；要么不得不采取种种狡辩以避免得罪天主，而这即使在军旅之外也不被容许；要么最终，为天主的缘故，必须忍受那公民信德同样已准备好接受的命运。军旅之役不能提供逃避罪罚或免于殉道的指望。基督徒在任何地方都不会改变其品格。只有一个福音，同一位耶稣，祂将来要否认凡否认祂的人，承认凡承认天主的人；祂也要救那为祂而丧失的生命，但另一方面，要毁灭那为获利而保全以致羞辱祂的生命。在祂眼中，忠信公民是士兵，正如忠信士兵是公民。信德状态不接受以\Quote{必要}为借口；他们唯一必要的是不犯罪，因此他们无必要犯罪。因为即使有人因酷刑或惩罚的必要而被强迫献祭并彻底否认基督，纪律也不纵容那种必要，因为恐惧否认基督和接受殉道的必要性高于逃避苦难和履行要求。事实上，这种借口推翻了我们圣事的整个本质，甚至消除了自愿犯罪的障碍；因为同样可以辩称倾向是必要的，因为它确实包含某种强制。事实上，我已经处理了这种关于必要的指控，它常被用来为与官职相关的冠冕辩护；因为正是为此，要么必须拒绝官职以免陷入罪恶行为，要么忍受殉道以摆脱官职。关于这问题的首要方面——甚至军事生活本身的非法性——我不再多加论述，以便次要问题回归其位。的确，若我全力排除军事生活，那么我现在就军冠问题发出挑战便毫无意义。那么，假定就冠冕的辩护而言，军役是合法的。\par

% structure: p0025 source=h2 target=section reason=relative-heading-rank; base=h2

\section{第十二章}

但我首先要就冠冕本身说几句话。这顶桂冠是奉献给\Person{阿波罗}或\Person{巴克斯}的——前者作为弓术之神，后者作为胜利之神。\Person{克劳狄}同样教导说，士兵们也常以桃金娘编成花冠。因为桃金娘属于\Person{维纳斯}，\OriginalTerm{埃涅阿斯后代}{Æneadæ}之母，亦是战神的情妇，她通过\Person{伊利亚}和\Person{罗慕路斯}而成为罗马的主保。但我不相信\Person{维纳斯}和\Person{玛尔斯}一样是罗马的，因为那情妇给她带来了烦恼。当军役又以橄榄冠为荣时，这偶像崇拜是向\Person{弥涅瓦}致敬，她同样是武器女神——但她得到了上述树木的冠冕，因为她与\Person{尼普顿}达成了和平。就此而言，军冠的迷信将处处受到玷污，且玷污一切。而且我认为，它在根源上也更加受玷污。看哪，每年公开宣发的誓愿，表面上看起来是什么？它们首先在军营中将军帐篷所在处举行，然后又在神庙中举行。除地点外，也请注意这些话语：\Quote{我们誓愿：\Person{朱庇特}啊，那时你将有一只双角镀金的公牛。}这话语是什么意思？无疑是指否认基督。尽管基督徒在这些地方口里不说一字，他头上戴冠便作了回应。在分发赏金时，同样命令使用桂冠。所以你看到，偶像崇拜并非无利可图，它出卖基督以换取金币，正如犹大以银币出卖他一样。难道\BibleQuote{你不能同时事奉天主而又事奉\OriginalTerm{玛门}{mammon}}\BibleRef{玛 6:24}这句话，意味着要全力以赴事奉\OriginalTerm{玛门}{mammon}，而背离天主吗？难道\BibleQuote{凯撒的归凯撒，天主的归天主}\BibleRef{玛 22:21}这句话，不仅是说不要把本人交还给天主，甚至还要从\Person{凯撒}那里取回银币吗？胜利的桂冠是由树叶编成的，还是由尸体编成的？它是以彩带装饰的，还是以坟墓装饰的？它是用香膏浸湿的，还是用妻子和母亲的泪水浸湿的？这或许也出自某些基督徒；因为基督也在外邦人中。那个为此原因（戴冠）在头上的人，岂不是甚至与自己作战吗？另一个服役之子属于皇家卫队。事实上，这些冠冕被称为\OriginalTerm{营冠}{Castrenses}，因其属于军营；同样，\Emph{\OriginalTerm{俸禄冠}{Munificæ}}也因其履行的凯撒职责而得名。但即使如此，你仍是别人的士兵和仆人；如果事奉两位主人，即天主和凯撒，那么当我以为，在天主具有更高要求的那些事上，甚至在两者都有利益的事上，你确实不属于凯撒。\par

% structure: p0027 source=h2 target=section reason=relative-heading-rank; base=h2

\section{第十三章}

为了国家的缘故，各等级公民也佩戴月桂花冠；但官员还佩戴金冠，如在雅典和罗马。甚至这些都比不上伊特鲁里亚冠冕。这一名称给予了那些以宝石和金橡叶为特色的冠冕，佩戴者穿着绣有棕榈枝的斗篷，引导载有神像的战车前往竞技场。还有行省的金冠，如今需要比人头更大的雕像头部来佩戴。但你们的等级、你们的官职，以及你们的集会场所——教会，都属于基督。你们属于祂，因为你们已被登记在生命册上。\BibleRef{斐 4:3} 在那里，主的血为你们作紫袍，祂的十字架为你们作宽饰带；在那里，斧头已经放在树根上；\BibleRef{玛 3:10} 那里有从叶瑟之根发出的枝条。\BibleRef{依 11:1} 不必在意国家的马匹及其冠冕。你们的主，当祂按圣经所述，凯旋进入耶路撒冷时，连一匹自己的驴都没有。这些人（依靠）战车，那些人（依靠）马匹；但我们要因上主我们天主的名而寻求帮助。我们甚至被召唤离开\BibleBook{若望}{默示录}中那个巴比伦的居所，更何况它的浮华。暴民也佩戴冠冕，或是因为皇帝胜利的盛大欢庆，或是因为市政节日的习俗。奢华试图将每一次公众欢乐的机会都据为己有。但至于你，你在这个世界是异乡人，是天上耶路撒冷的公民。宗徒说，我们的家乡在天上。\BibleRef{斐 3:20} 你有自己的登记册，自己的历法；你与世界的欢乐无关；相反，你被召唤去面对完全相反的事，因为\Quote{世界将欢乐，但你们将哀悼。}\BibleRef{若 16:20} 我想主宣告，哀恸的人是有福的，而非佩戴冠冕的人。婚姻也用冠冕装饰新郎；因此，我们不要异教的新娘，以免她们引诱我们陷入她们婚姻中伴随的偶像崇拜。你们从先祖那里得到了律法；你们有宗徒吩咐人在主内结婚。\BibleRef{格前 7:39} 在获释成为自由人时，也有加冕；但\Emph{你们}已经被基督赎回了，而且代价巨大。世界如何能释放另一个人的仆人？它看似自由，却终将成为束缚。在世上，一切都是名义上的，没有什么是真实的。因为即使当时，你们已被基督赎回，不受任何人的奴役；而现在，尽管人给了你们自由，你们却是基督的仆人。如果你们认为世上的自由是真实的，甚至用冠冕来确认它，那么你们就回到了人的奴役之下，误以为那是自由；你们失去了基督的自由，却误以为那是奴役。关于佩戴冠冕的原因，还有什么可争论的？竞技比赛轮流提供冠冕，而这些冠冕因其既献给神明又纪念死者，其本身的理由就足以定它们的罪。剩下的，就是\OriginalTerm{奥林匹亚的朱庇特}{Olympian Jupiter}、\OriginalTerm{内梅亚的赫拉克勒斯}{Nemean Hercules}、\OriginalTerm{阿尔克莫鲁斯}{Archemorus}、\OriginalTerm{安提诺乌斯}{Antinous}，都应该在一个基督徒身上加冕，使他本人成为令人厌恶的景观。我想，我们已经列举了所有佩戴冠冕的原因，没有一个与我们有关：它们都是与我们不相干的、不圣洁的、非法的，在圣事的庄严宣告中已经永远被弃绝。因为它们属于魔鬼及其使者的浮华，世上的官职、荣誉、节庆、猎取名声、虚假誓愿、人类奴性的表现、空洞的赞美、可耻的荣耀，其中都包含偶像崇拜，即使仅仅考虑冠冕的起源，一切冠冕都与此相连。\Person{克劳狄乌斯}确实会在他的序言中告诉我们，在\Person{荷马}的史诗中，天空也被星座加冕，而这无疑来自天主，无疑为了人；因此，人本身也应该被天主加冕。但世界给妓院、浴场、面包房、监狱、学校、甚至竞技场、剥去死去角斗士衣服的密室、死人的棺架加冕。这一服饰是多么神圣、圣洁、可敬和纯洁，不仅要从诗歌的天空，而且要从整个世界的交易中判断。然而，一个基督徒甚至不会用月桂花冠玷污自己的大门，如果他知道魔鬼将多少神明依附在门上：\OriginalTerm{雅努斯}{Janus}得名于门，\OriginalTerm{利门提努斯}{Limentinus}得名于门槛，\OriginalTerm{福尔库斯}{Forcus}和\OriginalTerm{卡尔纳}{Carna}得名于叶子和铰链；在希腊人中，还有\OriginalTerm{提瑞亚的阿波罗}{Thyræan Apollo}和恶灵\OriginalTerm{安特利伊}{Antelii}。\par

% structure: p0029 source=h2 target=section reason=relative-heading-rank; base=h2

\section{第十四章}

基督徒更不能将偶像崇拜的事奉戴在自己头上——不，我几乎可以说，是戴在基督头上，因为基督是基督徒的头——（因为他的头）是自由的，如同基督本身一样自由，没有义务佩戴遮盖物，更不必说束带\OriginalTerm{束带}{band}了。但是，那必须佩戴面纱\OriginalTerm{面纱}{veil}的头，我说的是女人的头，既然已经被这件事所占有，也不能再戴束带。她必须承担她自身谦卑的重担。如果她因天使的缘故不该不蒙头出现，那么她头上戴着冠冕\OriginalTerm{冠冕}{crown}岂不更冒犯那些（长老们），他们可能那时正戴着冠冕在天上？\BibleRef{默 4:4}因为女人的头上戴着冠冕，不是被诱人的美，而是完全放荡的标志——一个引人注目的廉耻丢弃，使诱惑燃烧起来吗？因此，女人要听从宗徒们的预见，不要过分地装饰自己，以免因任何精巧的发髻而加冕。但是，请问，那作为男人的头和女人的光荣的基督耶稣，教会的净配，为男女两性忍受了怎样的花环\OriginalTerm{花环}{garland}？我想，是荆棘\OriginalTerm{荆棘}{thorns}和蒺藜——这是肉身土壤为我们生出的罪过之象征，但十字架的能力移除了它们，借着我们主的头忍受了这一切，使死亡的每一根刺都变得迟钝。是的，除了象征之外，还有诋毁的嘴唇，羞辱，不名誉，以及当时凌辱和撕裂主之太阳穴的残酷事物中所含的凶暴，使你如今可以用月桂、桃金娘、橄榄和任何著名的枝条来加冕，更有用的是，还有从弥达斯的花园中采摘的百瓣玫瑰，以及两种百合花，各种紫罗兰，也许还有宝石和金子，甚至可以与基督后来获得的那冠冕相匹敌。因为他在尝了苦胆之后才尝了蜂蜜，他并非在天上被尊为光荣的君王，直到他被定罪为犹太人的君王，先在父面前暂时比天使微小一点，然后才被尊贵和光荣加冕。如果为了这些，你应将你的头归给他，若有可能，就请偿还吧，正如他将他的头替你呈现出来；否则，就不要以花朵加冕，因为你不能以荆棘加冕，因为你不可以以花朵加冕。\par

% structure: p0031 source=h2 target=section reason=relative-heading-rank; base=h2

\section{第十五章}

为天主保守祂自己的产业不受玷污；祂若愿意，必要给它加冕。不，祂甚至愿意。祂召唤我们。祂对得胜者说：\Quote{我要将生命的冠冕赐给你}。\BibleRef{弟后 4:8}你\Emph{你}也要至死忠心，你\Emph{你}也要打这场好仗，宗徒感觉如此确信那冠冕\OriginalTerm{冠冕}{crown}已为他存留。那骑白马出去得胜又要得胜的天使\BibleRef{默 6:2}得了胜利的冠冕；另一位\BibleRef{默 10:1}被一圈虹彩（如同其美丽的颜色——一片天上的草原）所装饰。同样，长老们座位周围戴着冠冕，也戴着金冠冕，而人子自己在云层之上闪耀。如果这些是在先知的异象中的显现，那么在真实显现中会是怎样的现实呢？看看那些冠冕。闻闻那些香气。为何你将那本应归于王冠\OriginalTerm{王冠}{diadem}的额角判给一个小小的花冠\OriginalTerm{花冠}{chaplet}或一条扭结的发带？因为基督耶稣已使我们成为天主和祂父的君王。你与那即将凋谢的花朵有何关系？你在叶瑟的枝条\OriginalTerm{叶瑟的枝条}{Branch of Jesse}上有一朵花，圣神的恩宠丰盛地栖息其上——一朵纯洁、不朽、永恒的花，选择这朵花的良好战士也在天军中得到晋升。羞愧吧，他的战友们，从此不要再被那战士定罪，而是被某个密特拉\OriginalTerm{密特拉}{Mithras}的士兵定罪，他在阴暗的洞穴中受培育时（可以说是黑暗的军营），当在刀剑下向他呈上一顶冠冕，如同殉道\OriginalTerm{殉道}{martyrdom}的滑稽模仿，随即放在他头上时，他得到告诫要抵抗并扔掉它，如果愿意，还可以把它转移到肩膀上，说密特拉就是他的冠冕。从此他再也不会被加冕；他以此为标志表明他是谁，如果他在任何地方因他的宗教受到考验，他如果扔掉冠冕，就被立即相信是密特拉的士兵——如果他说在他的神那里他有他的冠冕。让我们留意魔鬼的诡计，他惯于模仿天主的一些事物，别无他意，而是借着他仆人们的忠诚，来使我们蒙羞和定罪。\par

\SourceRecord{Translated by S. Thelwall.；Ante-Nicene Fathers；Vol. 3.；Edited by Alexander Roberts, James Donaldson, and A. Cleveland Coxe.；Buffalo, NY: Christian Literature Publishing Co.,；1885.；<http://www.newadvent.org/fathers/0304.htm>.}
